% Standalone problem document, generated from the adjacent README.md.
% Compile with XeLaTeX. Mathematical content is maintained in Markdown.
\documentclass[11pt,a4paper]{article}
\usepackage[fontsize=10.5pt]{scrextend}
\usepackage[margin=22mm,headheight=15pt,headsep=9mm,footskip=11mm]{geometry}
\usepackage{fontspec}
\setmainfont{texgyrepagella}[Extension=.otf,UprightFont=*-regular,BoldFont=*-bold,ItalicFont=*-italic,BoldItalicFont=*-bolditalic]
\setsansfont{texgyreheros}[Extension=.otf,UprightFont=*-regular,BoldFont=*-bold,ItalicFont=*-italic,BoldItalicFont=*-bolditalic]
\setmonofont{texgyrecursor}[Extension=.otf,UprightFont=*-regular,BoldFont=*-bold,ItalicFont=*-italic,BoldItalicFont=*-bolditalic,Scale=MatchLowercase]
\usepackage{amsmath,amssymb}
\usepackage{unicode-math}
\setmathfont{texgyrepagella-math.otf}
\usepackage{xcolor,microtype,enumitem,needspace,fancyhdr}
\usepackage{bookmark}
\definecolor{ink}{HTML}{17344B}
\definecolor{accent}{HTML}{197D80}
\definecolor{muted}{HTML}{596773}
\hypersetup{colorlinks=true,linkcolor=accent,urlcolor=accent,pdftitle={MI-28 - Ghabries's
determinant comparison for arbitrary nonnegative base
powers},pdfauthor={Open Problems in Numerical Linear Algebra}}
\urlstyle{same}
\setlength{\parindent}{0pt}
\setlength{\parskip}{5pt plus 1pt minus 1pt}
\linespread{1.04}
\setlength{\emergencystretch}{3em}
\widowpenalty=10000
\clubpenalty=10000
\displaywidowpenalty=10000
\allowdisplaybreaks[1]
\setlist{nosep,leftmargin=1.5em}
\providecommand{\tightlist}{\setlength{\itemsep}{0pt}\setlength{\parskip}{0pt}}
\providecommand{\pandocbounded}[1]{#1}
\pagestyle{fancy}
\fancyhf{}
\fancyhead[L]{\sffamily\footnotesize\color{muted}OPEN PROBLEMS IN NUMERICAL LINEAR ALGEBRA}
\fancyhead[R]{\sffamily\bfseries\footnotesize\color{accent}MI-28}
\fancyfoot[L]{\parbox[b]{0.9\textwidth}{\sffamily\fontsize{8}{10}\selectfont\color{muted}Editorial ratings; a bounded search cannot certify that a problem remains unsolved.\\Verification check: 2026-09-18}}
\fancyfoot[R]{\sffamily\footnotesize\color{muted}\thepage}
\renewcommand{\headrulewidth}{0.3pt}
\renewcommand{\headrule}{\hbox to\headwidth{\color{accent}\leaders\hrule height \headrulewidth\hfill}}
\makeatletter
\renewcommand{\section}{\Needspace{5\baselineskip}\@startsection{section}{1}{0pt}{12pt plus 2pt minus 2pt}{4pt}{\sffamily\bfseries\large\color{ink}}}
\renewcommand{\subsection}{\Needspace{4\baselineskip}\@startsection{subsection}{2}{0pt}{9pt plus 1pt minus 1pt}{3pt}{\sffamily\bfseries\normalsize\color{ink}}}
\makeatother
\setcounter{secnumdepth}{0}
\begin{document}
{\sffamily\small\color{muted}Matrix inequalities and norms\par}
\vspace{3pt}
{\raggedright\hyphenpenalty=10000\exhyphenpenalty=10000\sffamily\bfseries\fontsize{21}{24}\selectfont\color{ink}Ghabries's
determinant comparison for arbitrary nonnegative base powers\par}
\vspace{7pt}
{\sffamily\small\textbf{Difficulty:} challenging\quad\textbf{Importance:} interesting
to specialist\par}
{\sffamily\small\bfseries\color{accent}Status: Lean verified\par}
\vspace{4pt}
{\color{accent}\hrule height 0.8pt}
\vspace{6pt}
\textbf{Rating rationale (historical):} The intermediate base powers
resist the available determinantal comparisons, making the problem
challenging; the immediate application is to specialist matrix-function
inequalities.

\subsection{Resolution --- 2026-09-11}\label{resolution-2026-09-11}

\textbf{Solved (affirmative).} George Stepaniants's
\href{https://github.com/ajt60gaibb/OpenProblemsInNLA/blob/main/matrix-inequalities-and-norms/MI-28/solution.md}{complete
proof}, \textbf{Theorem 1 and Corollary 6}, proves the displayed
determinant inequality for every dimension, every complex positive
definite pair, all \(k\ge0\) and all \(0\le p\le2\). It proves the
stronger log-majorization of the corresponding normalized matrices. The
proof credits Ghabries--Abbas--Mourad--Assi's published result for
\(k\ge2\) and establishes the remaining range using Furuta inequalities
and a parameter interchange.
\href{https://github.com/ajt60gaibb/OpenProblemsInNLA/blob/main/matrix-inequalities-and-norms/MI-28/solution.pdf}{Proof
PDF} ·
\href{https://github.com/ajt60gaibb/OpenProblemsInNLA/blob/main/matrix-inequalities-and-norms/MI-28/solution.tex}{Standalone
TeX}.

The full analytic argument passed a
\href{https://github.com/ajt60gaibb/OpenProblemsInNLA/blob/main/references/stepaniants-mi28-2026-09-11/verification/reviews/MI-28-review.md}{separate
Codex-agent review}, including every exponent restriction, the
noncommuting factor order, exterior powers, endpoints and exact source
substitutions. It was developed with ChatGPT/Codex; that September 11
verification was independent agent review. The later Lean verification
is recorded below; no external human peer review is claimed.
\href{https://github.com/ajt60gaibb/OpenProblemsInNLA/blob/main/references/stepaniants-mi28-2026-09-11/README.md}{Authorship,
source checks, reproduction and public-branch audit}. The original ID,
path, target, historical ratings and earlier status check below are
retained.

\newpage

\subsection{Lean proof and verification
evidence}\label{lean-proof-and-verification-evidence}

\textbf{The complete original target is Lean verified, 2026-09-18
(UTC).} The
\href{https://github.com/sgstepaniants/OpenProblemsInNLA/blob/aae2941f9a7f1e28b8010f5216574d93bdc72181/matrix-inequalities-and-norms/MI-28/lean/Solution.lean}{immutable
proof} at revision \textit{aae2941f9a7f1e28b8010f5216574d93bdc72181}
passed
\href{https://github.com/sgstepaniants/OpenProblemsInNLA/actions/runs/35374928604/job/105697400237}{non-root
Linux run 35374928604}. The
\href{https://github.com/ajt60gaibb/OpenProblemsInNLA/blob/main/matrix-inequalities-and-norms/MI-28/lean/verification/linux-2026-09-18/README.md}{retained
execution evidence} binds all 282 submitted project inputs and 20
exported statements. Real Comparator, default-kernel replay, standard
transitive axioms, sandbox checks and rejection controls passed. Later
publication and PR merge checkouts are checked separately.

The main declaration \textit{NLA.MI28.determinant\_comparison} proves
the unchanged determinant inequality for every positive dimension, all
complex positive definite matrices, every real nonnegative base power
and the full closed power interval in the original statement.
\textit{NLA.MI28.full\_log\_majorization} proves the stronger normalized
comparison. Both determinant reality and positivity are proved,
including for the non-Hermitian right matrix. Powers and modulus use
concrete spectral functional calculus; no commutativity, spectral gap or
upper bound on the base power is assumed. Zero exponents, both interval
endpoints and every exterior degree are included. The
\href{https://github.com/ajt60gaibb/OpenProblemsInNLA/blob/main/matrix-inequalities-and-norms/MI-28/lean/Challenge.lean}{20
exact contracts} and
\href{https://github.com/ajt60gaibb/OpenProblemsInNLA/blob/main/matrix-inequalities-and-norms/MI-28/lean/PROOF-REVIEWER-MAP.md}{proof
map} identify all definitions and assumptions. Furuta and the large-base
comparison are proved internally, with prior mathematical credit
retained; no cited theorem is a Lean axiom.

Mathematical resolution and formalization: \textbf{George Stepaniants},
Department of Computing and Mathematical Sciences, California Institute
of Technology. The original question and prior results of Ghabries,
Abbas, Mourad and Assi, and Furuta, Fujii and Tanahashi background
retain their attribution. Two independent statement reviews preceded
implementation; two independent nonauthor final reviews approved the
full 52-module proof closure, including 18 unchanged, credited MI24
modules. The
\href{https://github.com/ajt60gaibb/OpenProblemsInNLA/blob/main/matrix-inequalities-and-norms/MI-28/lean/REVIEW-INDEX.md}{review
records} disclose substantial ChatGPT/Codex assistance and distinguish
source review from actual execution.

The project pins Lean 4.33.1,
\href{https://github.com/leanprover-community/mathlib4/tree/0df444a360eaa60ab8c11dca51a86af692955474}{Mathlib}
and
\href{https://github.com/alerad/leancert/tree/621a43d7cf21f87872392a01e874f2f1dbddc926}{LeanCert}.
Kernel-mode LeanCert certifies the fixed half-exponent used by the
modulus-square argument; the matrix and unbounded-parameter reasoning is
symbolic. All exports use only \textit{propext},
\textit{Classical.choice} and \textit{Quot.sound}. Run
\textit{lake build Solution} in
\textit{matrix-inequalities-and-norms/MI-28/lean};
\href{https://github.com/ajt60gaibb/OpenProblemsInNLA/blob/main/matrix-inequalities-and-norms/MI-28/lean/README.md}{reproduction
instructions} distinguish local compilation from the genuine Linux
checker.

\newpage

\subsection{Problem statement}\label{problem-statement}

For every integer \(n\ge1\), positive definite matrices
\(A,B\in\mathbb C^{n\times n}\), and real parameters \(k\ge0\),
\(0\le p\le2\), determine whether

\[\det(A^k+|AB|^p)\ge\det(A^k+A^pB^p)\]

holds. Here \(|AB|=((AB)^*(AB))^{1/2}\), and powers are defined by
spectral functional calculus; the zeroth power of a positive definite
matrix is the identity. Although the matrix inside the determinant on
the right need not be Hermitian, its determinant is a positive real
number: factor out \(A^k\) and use that a product of two positive
definite matrices has positive eigenvalues.

The source states the conjecture for positive semidefinite inputs. The
positive definite formulation captures the nontrivial question and
avoids ambiguity in zeroth powers; positive-exponent semidefinite cases
follow by regularization.

Such determinants arise in comparisons of interpolants for positive
definite data. The distinction between a product's modulus and the
product of matrix powers is important in matrix-function inequalities.

\subsection{References}\label{references}

\begin{enumerate}
\def\labelenumi{\arabic{enumi}.}
\tightlist
\item
  M. M. Ghabries, \emph{Contributions to Matrix Inequalities and Some
  Applications}, PhD thesis, University of Angers and Lebanese
  University (2022), §2.5 and final Open Problems, Problem 1,
  pp.~109--110.
  \href{https://www.researchgate.net/publication/361793582_Contributions_to_Matrix_Inequalities_and_Some_Applications}{Thesis
  uploaded by its author}.
\item
  H. Abbas, M. M. Ghabries and B. Mourad, \emph{New determinantal
  inequalities concerning Hermitian and positive semi-definite
  matrices}, Operators and Matrices 15 (2021), 105--116, §4, Conjecture
  2, p.~116 (earlier, narrower parameter region).
  \href{https://files.ele-math.com/articles/oam-15-07.pdf}{Publisher
  PDF}.
\item
  M. M. Ghabries et al., \emph{A proof of a conjectured determinantal
  inequality}, Linear Algebra Appl. 605 (2020), 21--28, main theorem.
  \href{https://doi.org/10.1016/j.laa.2020.07.013}{Publisher}.
\end{enumerate}

Status check (2026-09-10): the thesis retains the range \(0< k<2\),
\(0< p<2\), after recording proved results for \(k\ge2\) and for
\(p=2\). Its broader question supersedes the 2021 formulation. The 2020
resolution concerns Lin's earlier \(k=2\) question. The author's July
2026 normal-matrix determinant paper and searches for ``Ghabries
determinant arbitrary k conjecture'', exact titles, and 2025/2026
yielded no full resolution. This is a bounded check.
\end{document}
